% Thesis abstract (page numbering starts here)
\chapter*{TÓM TẮT ĐỒ ÁN}
\addcontentsline{toc}{chapter}{Tóm tắt đồ án}

Đồ án nghiên cứu bài toán dự báo xu hướng biến động giá cổ phiếu trên thị trường chứng
khoán Việt Nam bằng cách kết hợp tin tức tài chính tiếng Việt với dữ liệu giá theo chuỗi
thời gian. Bối cảnh của bài toán xuất phát từ thực tế giá cổ phiếu không chỉ phản ánh
lịch sử giao dịch mà còn phản ứng trước những thông tin về kết quả kinh doanh, kế hoạch
đầu tư, thay đổi quản trị, chính sách, rủi ro pháp lý và các sự kiện bất thường. Phạm
vi nghiên cứu tập trung vào sàn HOSE và một nhóm mã thuộc VN30. Mục tiêu là xác định
hướng biến động của phiên giao dịch kế tiếp theo ba lớp tăng, đi ngang và giảm, thay vì
dự báo trực tiếp giá trị tuyệt đối.

Hệ thống được tổ chức thành hai nhánh. Nhánh thứ nhất sử dụng PhoBERT, một mô hình
Transformer tiền huấn luyện cho tiếng Việt, để phân loại tin tức tài chính thành ba
lớp tích cực, trung tính và tiêu cực \cite{Nguyen2020_200300744v3}. Nhãn được hiểu theo
tác động kỳ vọng của nội dung đối với doanh nghiệp hoặc cổ phiếu liên quan, không chỉ
theo sắc thái chung của câu văn. Sau khi suy luận, hệ thống giữ lại xác suất của cả ba
lớp thay vì chỉ giữ nhãn có xác suất cao nhất. Các xác suất này được tổng hợp theo mã
cổ phiếu và ngày giao dịch cùng số lượng tin, tỷ lệ lớp, độ phân tán và cờ cho biết có
tin.

Nhánh thứ hai sử dụng dữ liệu OHLCV điều chỉnh cùng các đặc trưng nhân quả gồm lợi suất,
lợi suất log, tỷ lệ giá đóng cửa với đường trung bình động, thay đổi khối lượng và độ
biến động. Hai nhánh được căn chỉnh theo mã và ngày sau khi áp dụng quy tắc thời điểm.
Tin được chuyển về múi giờ Việt Nam; tin sau mốc cắt thông tin, tin cuối tuần hoặc tin
vào ngày nghỉ được ánh xạ sang phiên giao dịch kế tiếp. Tại ngày quan sát $t$, hệ thống
chỉ sử dụng thông tin có thể biết đến ngày đó, còn giá của phiên kế tiếp chỉ dùng để tạo
nhãn mục tiêu.

Về kiến trúc, các mô hình được xếp theo từng tầng để tách câu hỏi về dữ liệu khỏi câu
hỏi về độ phức tạp. Tầng đầu gồm mô hình luôn chọn lớp phổ biến nhất và mô hình ngẫu
nhiên theo phân bố lớp. Tầng cơ sở gồm LSTM chỉ sử dụng giá và LSTM hai nhánh kết hợp
giá với cảm xúc. Bộ biến đổi hợp nhất theo thời gian chỉ là phương án mở rộng khi dữ
liệu, tài nguyên và kết quả của cấu hình cơ sở cho phép.

Về phương pháp, đồ án xây dựng quy trình xử lý từ khảo sát tài liệu, thu thập và làm sạch
dữ liệu, loại bỏ bản ghi trùng, lưu thông tin truy nguyên, gắn tin với mã cổ phiếu, đồng
bộ theo phiên, tinh chỉnh mô hình cảm xúc, tạo đặc trưng, huấn luyện và đánh giá. Bộ dữ
liệu tin tức gồm tiêu đề, nội dung khi có thể thu thập, nguồn, URL, thời điểm công bố và
mã liên quan. Những bài đề cập nhiều mã được lưu theo quan hệ nhiều mã; những bản ghi
thiếu thời điểm hoặc không xác định được mã được đánh dấu và không đưa trực tiếp vào
bảng đặc trưng chính.

Dữ liệu được chia theo thứ tự thời gian bằng kiểm định cuốn chiếu (cuốn chiếu theo thời gian), có tập
kiểm tra cuối được khóa sau khi lựa chọn cấu hình. Bộ chuẩn hóa, bộ điền khuyết, ngưỡng
tạo nhãn, lựa chọn đặc trưng và các tham số khác chỉ được học từ phần huấn luyện của
từng cửa sổ. Mô hình có cảm xúc được so sánh với mô hình chỉ dùng giá trên cùng ngày
huấn luyện, xác thực và kiểm thử, cùng cách tiền xử lý, hạt giống ngẫu nhiên hoặc tập
hạt giống và tiêu chí dừng. Thiết kế này nhằm hạn chế rò rỉ dữ liệu, đồng thời tách phần
đóng góp của tín hiệu cảm xúc khỏi ảnh hưởng của kiến trúc hoặc cách chia mẫu.

Về dữ liệu, kho tin đã lưu 10.695 bản ghi có URL, thời điểm đăng và nội dung, kết hợp với
12.624 liên kết tin tức và mã cổ phiếu thuộc 10 mã VN30, cùng 14.990 dòng dữ liệu giá
OHLCV trong giai đoạn 2020--2025. Tập nhãn cảm xúc trong miền dữ liệu gồm 306 mẫu do một
người rà soát, bên cạnh tập hạt giống CafeF gồm 999 mẫu.

Về kết quả, mô hình cảm xúc PhoBERT đạt F1 vĩ mô 0,5176 và độ chính xác cân bằng 0,5473
trên xác thực chéo phân tầng 5 lượt. Lớp tiêu cực chỉ có 31 mẫu và độ bao phủ 0,0619, nên lớp này
chưa được đo một cách đáng tin cậy. Ở bài toán dự báo, mô hình LSTM chỉ dùng giá đạt F1 vĩ
mô 0,3797, độ chính xác cân bằng 0,3915 và AUC-ROC một-chống-còn-lại 0,5721, vượt rõ mốc
ngẫu nhiên 0,3289 và mốc lớp phổ biến 0,1603. Vì nhãn ba lớp được cắt theo phân vị một
phần ba, các lớp gần cân bằng theo thiết kế và mức ngẫu nhiên của bài toán là 0,333.

Đóng góp riêng của tín hiệu cảm xúc được đo bằng một lần chạy đối chứng giữ nguyên kiến
trúc hai nhánh nhưng thay xác suất cảm xúc thật bằng phân phối trung tính hằng số. Phép
tách này cho thấy hiệu số trực tiếp $+0{,}0028$ F1 vĩ mô giữa mô hình hai nhánh và mô hình
chỉ dùng giá hoàn toàn đến từ hiệu ứng kiến trúc $(+0{,}0051)$; đóng góp thông tin thực của
cảm xúc là $-0{,}0024$ với khoảng tin cậy 95\% lấy mẫu lặp theo khối ngày là
$[-0{,}0113; +0{,}0111]$. Phân tầng theo cờ có tin cho hiệu số $-0{,}0009$ trên riêng nhóm
phiên thực sự có tin, nên giả thuyết tín hiệu bị pha loãng bởi các phiên không có tin bị
loại trừ.

Đồ án không đặt trước giả định rằng tín hiệu cảm xúc chắc chắn cải thiện dự báo. Kết luận
là chưa có bằng chứng về cải thiện trong giao thức này, đồng thời xác định được biên độ của
đóng góp đó nhỏ hơn khoảng $\pm0{,}011$ F1 vĩ mô. Một giới hạn quan trọng là chất lượng bộ
sinh cảm xúc chưa tách được khỏi độ mạnh của bản thân tín hiệu, nên đồ án không quy kết
nguyên nhân cho kết luận âm này. Kết quả không được diễn giải thành quan hệ nhân quả hay
khuyến nghị đầu tư.

\vspace{0.5cm}
\noindent Từ khóa: dự báo xu hướng giá cổ phiếu, phân tích cảm xúc, PhoBERT,
Transformer, chuỗi thời gian, HOSE, VN30.
